% Options for packages loaded elsewhere
% Options for packages loaded elsewhere
\PassOptionsToPackage{unicode}{hyperref}
\PassOptionsToPackage{hyphens}{url}
\PassOptionsToPackage{dvipsnames,svgnames,x11names}{xcolor}
%
\documentclass[
  ngerman,
  letterpaper,
  DIV=11]{scrreprt}
\usepackage{xcolor}
\usepackage{amsmath,amssymb}
\setcounter{secnumdepth}{5}
\usepackage{iftex}
\ifPDFTeX
  \usepackage[T1]{fontenc}
  \usepackage[utf8]{inputenc}
  \usepackage{textcomp} % provide euro and other symbols
\else % if luatex or xetex
  \usepackage{unicode-math} % this also loads fontspec
  \defaultfontfeatures{Scale=MatchLowercase}
  \defaultfontfeatures[\rmfamily]{Ligatures=TeX,Scale=1}
\fi
\usepackage{lmodern}
\ifPDFTeX\else
  % xetex/luatex font selection
\fi
% Use upquote if available, for straight quotes in verbatim environments
\IfFileExists{upquote.sty}{\usepackage{upquote}}{}
\IfFileExists{microtype.sty}{% use microtype if available
  \usepackage[]{microtype}
  \UseMicrotypeSet[protrusion]{basicmath} % disable protrusion for tt fonts
}{}
\makeatletter
\@ifundefined{KOMAClassName}{% if non-KOMA class
  \IfFileExists{parskip.sty}{%
    \usepackage{parskip}
  }{% else
    \setlength{\parindent}{0pt}
    \setlength{\parskip}{6pt plus 2pt minus 1pt}}
}{% if KOMA class
  \KOMAoptions{parskip=half}}
\makeatother
% Make \paragraph and \subparagraph free-standing
\makeatletter
\ifx\paragraph\undefined\else
  \let\oldparagraph\paragraph
  \renewcommand{\paragraph}{
    \@ifstar
      \xxxParagraphStar
      \xxxParagraphNoStar
  }
  \newcommand{\xxxParagraphStar}[1]{\oldparagraph*{#1}\mbox{}}
  \newcommand{\xxxParagraphNoStar}[1]{\oldparagraph{#1}\mbox{}}
\fi
\ifx\subparagraph\undefined\else
  \let\oldsubparagraph\subparagraph
  \renewcommand{\subparagraph}{
    \@ifstar
      \xxxSubParagraphStar
      \xxxSubParagraphNoStar
  }
  \newcommand{\xxxSubParagraphStar}[1]{\oldsubparagraph*{#1}\mbox{}}
  \newcommand{\xxxSubParagraphNoStar}[1]{\oldsubparagraph{#1}\mbox{}}
\fi
\makeatother


\usepackage{longtable,booktabs,array}
\usepackage{calc} % for calculating minipage widths
% Correct order of tables after \paragraph or \subparagraph
\usepackage{etoolbox}
\makeatletter
\patchcmd\longtable{\par}{\if@noskipsec\mbox{}\fi\par}{}{}
\makeatother
% Allow footnotes in longtable head/foot
\IfFileExists{footnotehyper.sty}{\usepackage{footnotehyper}}{\usepackage{footnote}}
\makesavenoteenv{longtable}
\usepackage{graphicx}
\makeatletter
\newsavebox\pandoc@box
\newcommand*\pandocbounded[1]{% scales image to fit in text height/width
  \sbox\pandoc@box{#1}%
  \Gscale@div\@tempa{\textheight}{\dimexpr\ht\pandoc@box+\dp\pandoc@box\relax}%
  \Gscale@div\@tempb{\linewidth}{\wd\pandoc@box}%
  \ifdim\@tempb\p@<\@tempa\p@\let\@tempa\@tempb\fi% select the smaller of both
  \ifdim\@tempa\p@<\p@\scalebox{\@tempa}{\usebox\pandoc@box}%
  \else\usebox{\pandoc@box}%
  \fi%
}
% Set default figure placement to htbp
\def\fps@figure{htbp}
\makeatother



\ifLuaTeX
\usepackage[bidi=basic]{babel}
\else
\usepackage[bidi=default]{babel}
\fi
% get rid of language-specific shorthands (see #6817):
\let\LanguageShortHands\languageshorthands
\def\languageshorthands#1{}
\ifLuaTeX
  \usepackage[german]{selnolig} % disable illegal ligatures
\fi


\setlength{\emergencystretch}{3em} % prevent overfull lines

\providecommand{\tightlist}{%
  \setlength{\itemsep}{0pt}\setlength{\parskip}{0pt}}



 


\KOMAoption{captions}{tableheading}
\makeatletter
\@ifpackageloaded{bookmark}{}{\usepackage{bookmark}}
\makeatother
\makeatletter
\@ifpackageloaded{caption}{}{\usepackage{caption}}
\AtBeginDocument{%
\ifdefined\contentsname
  \renewcommand*\contentsname{Inhaltsverzeichnis}
\else
  \newcommand\contentsname{Inhaltsverzeichnis}
\fi
\ifdefined\listfigurename
  \renewcommand*\listfigurename{Abbildungsverzeichnis}
\else
  \newcommand\listfigurename{Abbildungsverzeichnis}
\fi
\ifdefined\listtablename
  \renewcommand*\listtablename{Tabellenverzeichnis}
\else
  \newcommand\listtablename{Tabellenverzeichnis}
\fi
\ifdefined\figurename
  \renewcommand*\figurename{Abbildung}
\else
  \newcommand\figurename{Abbildung}
\fi
\ifdefined\tablename
  \renewcommand*\tablename{Tabelle}
\else
  \newcommand\tablename{Tabelle}
\fi
}
\@ifpackageloaded{float}{}{\usepackage{float}}
\floatstyle{ruled}
\@ifundefined{c@chapter}{\newfloat{codelisting}{h}{lop}}{\newfloat{codelisting}{h}{lop}[chapter]}
\floatname{codelisting}{Listing}
\newcommand*\listoflistings{\listof{codelisting}{Listingverzeichnis}}
\makeatother
\makeatletter
\makeatother
\makeatletter
\@ifpackageloaded{caption}{}{\usepackage{caption}}
\@ifpackageloaded{subcaption}{}{\usepackage{subcaption}}
\makeatother
\usepackage{bookmark}
\IfFileExists{xurl.sty}{\usepackage{xurl}}{} % add URL line breaks if available
\urlstyle{same}
\hypersetup{
  pdftitle={Zeitreihenanalyse},
  pdfauthor={Lukas Arnold; Simone Arnold; Florian Bagemihl; Matthias Baitsch; Marc Fehr; Maik Poetzsch; Sebastian Seipel},
  pdflang={de},
  colorlinks=true,
  linkcolor={blue},
  filecolor={Maroon},
  citecolor={Blue},
  urlcolor={Blue},
  pdfcreator={LaTeX via pandoc}}


\title{Zeitreihenanalyse}
\author{Lukas Arnold \and Simone Arnold \and Florian
Bagemihl \and Matthias Baitsch \and Marc Fehr \and Maik
Poetzsch \and Sebastian Seipel}
\date{2026-02-27}
\begin{document}
\maketitle

\renewcommand*\contentsname{Inhaltsverzeichnis}
{
\hypersetup{linkcolor=}
\setcounter{tocdepth}{0}
\tableofcontents
}

\bookmarksetup{startatroot}

\chapter*{Preamble}\label{preamble}
\addcontentsline{toc}{chapter}{Preamble}

\markboth{Preamble}{Preamble}

\phantomsection\label{Lizenz}
\begin{figure}

\begin{minipage}{0.20\linewidth}
\pandocbounded{\includegraphics[keepaspectratio]{index_files/mediabag/by.png}}\end{minipage}%
%
\begin{minipage}{0.80\linewidth}
Bausteine Computergestützter Datenanalyse. ``Zeitreihenanalyse'' von
Lukas Arnold, Simone Arnold, Florian Bagemihl, Matthias Baitsch, Marc
Fehr, Maik Poetzsch und Sebastian Seipel ist lizensiert unter
\href{https://creativecommons.org/licenses/by/4.0/deed.de}{CC BY 4.0}.
Das Werk ist abrufbar unter
\url{https://github.com/bausteine-der-datenanalyse/a-zeitreihenanalyse}.
Ausgenommen von der Lizenz sind alle Logos und anders gekennzeichneten
Inhalte. 2024\end{minipage}%

\end{figure}%

Zitiervorschlag

Arnold, Lukas, Simone Arnold, Matthias Baitsch, Marc Fehr, Maik
Poetzsch, und Sebastian Seipel. 2024. „Bausteine Computergestützter
Datenanalyse. Methodenbaustein Numerik``.
\url{https://github.com/bausteine-der-datenanalyse/a-zeitreihenanalyse}.

BibTeX-Vorlage

\begin{verbatim}
@misc{BCD-Styleguide-2024,
 title={Bausteine Computergestützter Datenanalyse. Methodenbaustein Numerik},
 author={Arnold, Lukas and Arnold, Simone and Baitsch, Matthias and Fehr, Marc and Poetzsch, Maik and Seipel, Sebastian},
 year={2024},
 url={https://github.com/bausteine-der-datenanalyse/a-zeitreihenanalyse}} 
\end{verbatim}

\bookmarksetup{startatroot}

\chapter*{Intro}\label{intro}
\addcontentsline{toc}{chapter}{Intro}

\markboth{Intro}{Intro}

\section*{Voraussetzungen}\label{voraussetzungen}
\addcontentsline{toc}{section}{Voraussetzungen}

\markright{Voraussetzungen}

\begin{itemize}
\tightlist
\item
  Python-Grundlagen
\item
  NumPy (Arrays, grundlegende Operationen)
\item
  Matplotlib (Visualisierung)
\item
  Grundverständnis von Pandas (Series, DataFrame)
\end{itemize}

\section*{Verwendete Pakete und
Datensätze}\label{verwendete-pakete-und-datensuxe4tze}
\addcontentsline{toc}{section}{Verwendete Pakete und Datensätze}

\markright{Verwendete Pakete und Datensätze}

\subsection*{Pakete}\label{pakete}
\addcontentsline{toc}{subsection}{Pakete}

\begin{itemize}
\tightlist
\item
  \href{https://numpy.org/}{NumPy}
\item
  \href{https://matplotlib.org/}{Matplotlib}
\item
  \href{https://pandas.pydata.org/}{Pandas}
\end{itemize}

\subsection*{Datensätze}\label{datensuxe4tze}
\addcontentsline{toc}{subsection}{Datensätze}

\section*{Bearbeitungszeit}\label{bearbeitungszeit}
\addcontentsline{toc}{section}{Bearbeitungszeit}

\markright{Bearbeitungszeit}

Geschätzte Bearbeitungszeit: 2h

\section*{Lernziele}\label{lernziele}
\addcontentsline{toc}{section}{Lernziele}

\markright{Lernziele}

\begin{itemize}
\tightlist
\item
  Zeitreihen strukturiert vorverarbeiten (Zeitindex, Missing, Ausreißer)
\item
  Resampling und Alignment methodisch korrekt einsetzen
\item
  Glättung so wählen, dass sie zur Fragestellung passt
\item
  Trend und Saisonanteile erkennen und interpretieren
\item
  zeitliche Abhängigkeiten mit Lag/Autokorrelation diagnostizieren
\item
  Anomalien mit einfachen, begründbaren Regeln detektieren
\item
  Forecast-Baselines implementieren und mit Fehlermaßen bewerten
\end{itemize}




\end{document}
